\begin{definition}
    Seien $\Aa, x \in \Aa, \Ba$ wie in Definition 1.27. Dann definieren wir
    \begin{equation*}
        \phi_x^t : C_b(\sigma(x), \C) \to C_b(M_\Ba, \C), f \mapsto f \circ \phi_x
    \end{equation*}
\end{definition}

\begin{lemma}
    Die Abbildung $\phi_x^t$ ist ein isometrischer $\ast$-Isomorphismus.
\end{lemma}

\begin{remark}
    Es gilt
    $$ \phi_x^t(\iota_{\sigma(x)}) = \iota_{\sigma(x)} \circ \phi_x = \hat{x} $$
    und
    $$ \phi_x^t(\zeta \mapsto \zeta^j \zeta^{-k}) = \hat{x}^j \hat{x}^{-k} = \widehat{x^j (x^*)^k}. $$
\end{remark}

\begin{definition}
    Sei $\Aa$ eine $C^*$-Algebra mit Eins, $x \in \Aa$ normal und $f \in C_b(\sigma(x), \C)$. Dann definieren wir
    $$ f(x) \coloneqq (\hat{\cdot})^{-1} \circ \phi_x^t(f). $$
\end{definition}

\begin{remark}{\ }
    \begin{enumerate}
        \item $(f \mapsto f(x)) = (\hat{\cdot})^{-1} \circ \phi_x^t(f)$ ist ein isometrischer $\ast$-Isomorphismus.
        
        Ist $f = \iota_{\sigma(x)}$, so gilt $f(x) = x$.
        
        Entsprechend gilt auch mit $f = (\zeta \mapsto \zeta^j \overline{\zeta}^k)$, dass $f(x) = x^j (x^*)^k$.

        \item $f$ ist reellwertig genau dann, wenn $f(x) = f(x)^*$.
        
        $f$ ist $\T$-wertig genau dann, wenn $f(x)$ unitär ist.

        $f(x) \in [e, x] \subseteq \Aa$.

        \item Ist $f \in C_b(\sigma(x), \C), a \in \Aa, ax = xa, ax^* = x^* a$, so ist auch $af(x) = f(x) a$.
        
        \item Ist $f \in C_b(\sigma(x), \C)$, so ist $f(x)$ das eindeutige Element von $[e, x]$ mit der Eigenschaft $\widehat{f(x)} = \phi_x^t(f)$, beziehungsweise $\hat{f}(x)(m) = \phi_x^t(f)(m)$ für alle $m \in M_\Ba$. Inbesondere gilt
        $$ \phi_x^t(f)(m) = (f \circ \phi_x)(m) = f(m(x)). $$

        \item Es gilt
        $$ \sigma_\Ba (f(x)) = \{ m(f(x)) : m \in M_\Ba \} = f(\{ m(x) : m \in M_\Ba \}) = f(\sigma(x)). $$
    \end{enumerate}
\end{remark}

\begin{proposition}
    Sei $\Aa$ eine $C^*$-Algebra mit Eins und $x \in \Aa$ normal. Sind $f \in C_b(\sigma(x), \C), g \in C_b(\sigma(f(x)), \C)$. Dann ist $g \circ f \in C_b(\sigma(x), \C)$ und $(g \circ f)(x) = g(f(x))$.
\end{proposition}

\begin{proof}
    Es gilt $f(x) \in [e, x]$ und damit $[e, f(x)] \leq [e, x] \leq \Aa$. Sei nun $m \in M_{[e,x]}$.
    \begin{align*}
        m(g(f(x))) &= \restr{m}{[e,f(x)]}{(g(f(x)))} = g(m(f(x))) = g(f(m(x))) \\
        &= (g \circ f)(m(x)) = m((g \circ f)(x)). \qedhere
    \end{align*}
\end{proof}

\begin{proposition}
    Seien $\Aa, \Ba$ $C^*$-Algebren mit Eins ($e_\Aa$, bzw. $e_\Ba$) und $\varphi : \Aa \to \Ba$ ein $\ast$-Algebren-Homomorphismus mit $\varphi(e_\Aa) = e_\Ba$. Falls $x \in \Aa$ normal ist, so ist auch $\varphi(x) \in \Ba$ normal mit $\sigma_\Ba(\varphi(x)) \subseteq \sigma_\Aa(x)$ und $\varphi(f(x)) = \restr{f}{\sigma_\Ba(\varphi(x))} (\varphi(x))$ für alle $f \in C_b(\sigma_\Aa(x), \C)$.\footnote{In anderen Worten: Der oben konstruierte Funktionalkalkül ist verträglich mit $\ast$-Algebren-Homomorphismen.}
\end{proposition}

\begin{proof}
    Wir wissen bereits $\Vert \varphi \Vert = 1$, $\varphi(x) \varphi(x)^* = \varphi(x)^* \varphi(x)$ und $\sigma_\Ba(\varphi(x)) \subseteq \sigma_\Aa(x)$.

    Sei $m \in M_{[e_\Ba, \varphi(x)]}$. Dann ist $m \circ \varphi \in M_{[e_\Aa, x]}$. Weiters gilt
    \begin{align*}
        m(\varphi(f(x))) &= (m \circ \varphi)(f(x)) = f((m \circ \varphi)(x)) = \restr{f}{\sigma_\Ba (\varphi(x))} (m(\varphi(x))) \\
        &= m(\restr{f}{\sigma_\Ba (\varphi(x))} (\varphi(x))). \qedhere
    \end{align*}
\end{proof}

Proposition 1.34 folgt aus Proposition 1.35 mit $f = \iota$.

\begin{corollary}
    Mit den Voraussetzungen von Proposition 1.35 gilt, dass wenn $\varphi$ injektiv ist, $\varphi$ sogar isometrisch ist.
\end{corollary}

\begin{proof}
    Angenommen $\varphi$ ist nicht isometrisch. Dann gibt es ein $a \in \Aa$ mit $\Vert a \Vert > \Vert \varphi(a) \Vert$, da $\Vert \varphi \Vert \leq 1$. Dann gilt
    $$ r(\varphi(a^* a)) = \Vert \varphi(a^* a) \Vert = \Vert \varphi(a)^* \varphi(a) \Vert = \Vert \varphi(a) \Vert^2 < \Vert a \Vert^2 = r(a^* a). $$
    Inbesondere folgt $\sigma_\Ba (\varphi(a^* a)) \subsetneq \sigma_\Aa(a^* a)$. Sei nun $f : \sigma(a^* a) \to \C$ stetig mit $f \neq 0$ und $\restr{f}{\sigma_\Ba (\varphi(a^* a))} = 0$. Damit erhalten wir aus der vorigen Proposition
    $$ \varphi(f(a^* a)) = \restr{f}{\sigma_\Ba(\varphi(a^* a))} (\varphi(a^* a)) = 0. $$
    Nun ist aber $f(a^* a) \neq 0$, daher $\ker \varphi \neq \{ 0 \}$ und damit $\varphi$ nicht injektiv.
\end{proof}

\begin{definition}
    Sei $\Aa$ eine $C^*$-Algebra (mit Eins) und $a \in \Aa$. Dann heißt $a$
    \begin{itemize}
        \item \emph{selbstadjungiert}, wenn $a^* = a$.
        \item \emph{positiv}, wenn $a^* = a$ mit $\sigma(a) \subseteq [0, +\infty)$.
    \end{itemize}
    Die Menge der selbstadjungierten, bzw. positiven Elemente von $\Aa$ bezeichnen wir mit $\Aa_{sa}$, bzw. $\Aa_+$.
\end{definition}

\begin{remark}
    Seien $a \in \Aa_{sa}$ und $p, q \in \Aa_+$.
    \begin{enumerate}
        \item Ist $\sigma(a) \in [0, +\infty)$, dann ist $\restr{\sqrt{\cdot}}{\sigma(a)} \in C_b(\sigma(a), \C)$. Insbesondere gilt
        $$ \sigma(\sqrt{q}) = \sqrt{\sigma(q)} \subseteq [0, +\infty). $$
        Weiters ist
        $$ (\sqrt{q})^2 = \sqrt{q} \sqrt{q} = (\sqrt{\cdot})^2(q) = q. $$
        Entsprechend sieht man $\sqrt{p^2} = p$.

        \item $\sqrt{q}$ ist das eindeutige Element aus $\Aa_+$ mit $(\sqrt{q})^2 = q$, wäre nämlich auch $p^2 = q = (\sqrt{q})^2$, so folgt $p = \sqrt{q}$.
        
        \item Für alle $n \in \N_0$ ist $q^n \in \Aa_+$ und $a^{2n} \in \Aa_+$.
        
        \item Betrachte die (stetige) Funktion
        $$ \vert \cdot \vert : \sigma(a) \to \R \ (\C), $$
        so ist $\vert a \vert \in \Aa_+$ (sogar wenn $a$ nicht selbstadjungiert ist, man beachte $\sigma(\vert a \vert) = \vert \sigma(a) \vert$). Weiters kann man die (stetigen) Funktionen
        $$ \max(0, \cdot) : \sigma(a) \to [0, +\infty), \quad -\min(0, \cdot) : \sigma(a) \to [0,+\infty), $$
        betrachten, und definieren
        $$ a_+ \coloneqq \max(0, \cdot)(a), \quad a_- \coloneqq [-\min(0,\cdot)](a), $$
        was beides in $\Aa_+$ liegt.

        \item Es gilt $\iota_{\sigma(a)} \in C_b(\sigma(a), \C)$ und damit $\Vert \iota_{\sigma(a)} \Vert = \Vert a \Vert$. Weiters gilt $0 \leq \max(0, \cdot) \leq \vert \iota_{\sigma(a)} \vert$ und damit
        $$ \Vert \max(0, \cdot) \Vert_\infty \leq \Vert \vert \iota_{\sigma(a)} \vert \Vert_\infty = \Vert a \Vert. $$
        Entsprechendes folgt für $-\min(0, \cdot)$ und damit
        $$ \Vert a_+ \Vert \leq \Vert a \Vert, \quad \Vert a_- \Vert \leq \Vert a \Vert, \quad \Vert \vert a \vert \Vert = \Vert a \Vert. $$

        Es gilt $\iota_{\sigma(a)} = \max(0, \cdot) - (- \min(0, \cdot))$, womit wir aus unserem $\ast$-Isomorphismus erhalten $ a = a_+ - a_- $. Entsprechend erhalten wir $\vert a \vert = a_+ + a_-$.

        Weiters gilt $0 = (\max(0, \cdot)) \cdot (-\min(0, \cdot))$ und damit $ a_+ a_- = 0 $.

        Außerdem gilt $\iota_{\sigma(a)}^2 = \vert \iota_{\sigma(a)} \vert^2$, damit $\vert a \vert^2 = a^2$ und insbesondere $\vert a \vert = \sqrt{a^2}$.

        \item Es gilt $\Aa_{sa} = \Aa_+ - \Aa_+$. Ähnlich leitet man her
        $$ \Aa = \Aa_{sa} + i \Aa_{sa} = \Aa_+ - \Aa_+ + i \Aa_+ - i \Aa_+. $$

        \item Sei $t \geq \Vert a \Vert$. Dann ist
        $$ a \in \Aa_+ \quad\Longleftrightarrow\quad \Vert te - a \Vert \leq t. $$
        Um dies einzusehen bemerken wir
        \begin{align*}
            \sigma(te - a) = t - \sigma(a) \subseteq [0, +\infty),
        \end{align*}
        sowie
        \begin{align*}
            \Vert te - a \Vert = r(te - a) = \max_{s \in \sigma(a)} (t-s) = t - \min_{s \in \sigma(a)} s.
        \end{align*}
        Also ist
        \begin{align*}
            a \in \Aa_+ \Leftrightarrow \min_{s \in \sigma(a)} s \geq 0 \Leftrightarrow \Vert te - a \Vert \leq t.
        \end{align*}

        \item Es ist zunächst $p + q \in \Aa_{sa}$. Setze $t \coloneqq \Vert p \Vert + \Vert q \Vert \geq \Vert p+q \Vert$, so ist
        $$ \Vert (\Vert p \Vert + \Vert q \Vert) e - (p+q) \Vert = \Vert \Vert p \Vert e - p \Vert + \Vert \Vert q \Vert e - q \Vert \leq \Vert p \Vert + \Vert q \Vert, $$
        womit aus dem vorigen Punkt sogar $p+q \in \Aa_+$ folgt.

        \item Seien $q_1, q_2 \in \Aa_+$ mit $q_1 q_2 = 0$ und $q_1 - q_2 = a \in \Aa_{sa}$. Dann gilt
        \begin{align*}
            (q_1 + q_2)^2 = q_1^2 + q_2^2 = (q_1 - q_2)^2 = a^2.
        \end{align*}
        Damit folgt $q_1 + q_2 = \vert a \vert$. Inbesondere folgt
        $$ q_1 = \frac{q_1 + q_2 + (q_1 - q_2)}{2} = \frac{\vert a \vert + a}{2} = a_+, $$
        $$ q_2 = \frac{q_1 + q_2 - (q_1 - q_2)}{2} = \frac{\vert a \vert - a}{2} = a_-. $$

        \item Wegen der Stetigkeit der $\ast$-Operation ist $\Aa_{sa} \subseteq \Aa$ abgeschlossen. Tatsächlich ist sogar $\Aa_+ \subseteq \Aa$ abgeschlossen. Um dies einzusehen sei $b_n$ eine Folge in $\Aa_+$ mit $b_n \to b \in \Aa_{sa}$. Dann gilt
        $$ \Vert \Vert b \Vert e - b \Vert = \lim_{n \to \infty} \Vert \Vert b_n \Vert e - b_n \Vert \leq \lim_{n \to \infty} \Vert b_n \Vert = \Vert b \Vert, $$
        und nach 7. folgt $b \in \Aa_+$.
    \end{enumerate}
\end{remark}

\begin{example}
    Sei $\Aa = L_b(H)$ für einen Hilbertraum $H$. Dann ist $T \in \Aa_{sa}$ genau dann wenn
    $$ \langle Tx, y \rangle = \langle x, Ty \rangle, $$
    und $T \in \Aa_+$ genau dann wenn
    $$ \langle Tx, x \rangle \geq 0. $$
    Inbesondere folgt (mit dem numerischen Wertebereich), dass aus positiv bereits selbstadjungiert folgt. Man kann also hier positiv, bzw. selbstadjungiert, über Elemente des Hilbertraums charakterisieren.
\end{example}